\subsection{$\SL_{2}$ 中的例子}
由定理 \ref{th:orb coho}，为计算轨道积分，不仅要知道仿射 Springer 纤维 $\sX_{\gamma}$ 的几何，还要知道 Frobenius 在其上同调上的作用。\S\ref{ss:aff examples} 中已见过代数闭域上仿射 Springer 纤维的许多例子，这里重点放在 Frobenius 作用上。

本小节取 $G=\SL_{2}$，并假设 $\chark>2$。我们在此情形计算几个轨道积分，并用显式计算验证定理 \ref{th:orb coho}。

\subsubsection{非分歧情形：$\gamma$}\label{unram sl2} 设 $\gamma=\mat{0}{at}{t}{0}\in\sl_{2}(F)$ 是正则半单元，其中 $a\in k^{\times}-(k^{\times})^{2}$。

设 $E\subset\mathfrak{gl}_{2}(F)$ 是 $\gamma$ 在 $\mathfrak{gl}_{2}(F)$ 中的中心化子。则 $E=\{u+v\gamma|u,v\in F\}$ 是 $F$ 的非分歧二次扩张，由添加 $\sqrt{a}$ 得到。因此 $E=k_{E}\lr{t}$，其中 $k_{E}= k(\sqrt{a})$。有 $G_{\gamma}(F)=(E^{\times})^{\Nm=1}=\ker(\Nm: E^{\times}\to F^{\times})$，它是紧群。固定 $G_{\gamma}(F)$ 上总体积为 $1$ 的 Haar 测度。

设 $\sX_{\gamma}$ 是 $\gamma$ 的仿射 Springer 纤维，它是 $k$ 上概形。由引理 \ref{l:Ogamma}，
\begin{equation*}
O_{\gamma}(\one_{\frg(\calO_{F})})=\#X_{\gamma}=\#\sX_{\gamma}(k).
\end{equation*}
在 \S\ref{P1 chain} 中已证明 $\sX_{\gamma,\kbar}$ 是可数多条有理曲线 $C_{n}\cong\PP^{1}$ 之并，指标为整数 $n\in\ZZ$。由于 $\gamma$ 在 $E$ 上可对角化，每个分支 $C_{n}$ 实际上定义在 $k_{E}=k(\sqrt{a})$ 上。格 $\L_{\gamma}\subset G_{\gamma}(F^{\ur})$ 包含在 $G_{\gamma}(E)\cong E^{\times}$ 中，并由素元 $t\in E$ 生成。把分支 $C_{n}$ 编号，使 $t\in \L_{\gamma}$ 把 $C_{n}$ 送到 $C_{n+1}$。令 $x_{n+1/2}:=C_{n}\cap C_{n+1}$，它是 $\sX_{\gamma}$ 的 $k_{E}$-点。

非平凡对合 $\sigma\in\Gal(k_{E}/k)$ 在 $G_{\gamma}(E)$ 上的作用是取逆，因此它在 $\L_{\gamma}$ 上的作用也是取逆。标准格 $\calO^{2}_{F}$ 同时落在 $C_{0}$ 与 $C_{1}$ 中，故它就是点 $x_{1/2}$。由于 $x_{1/2}$ 定义在 $k$ 上，它被 $\sigma$ 固定。$\sigma$ 在 $\sX_{\gamma,k_{E}}$ 上的作用与其在 $\L_{\gamma}$ 上的作用（取逆）相容，因此唯一可能是
\begin{equation*}
\sigma(C_{n})=C_{1-n}, \quad \sigma(x_{n+1/2})=x_{-n+1/2}, \quad\forall n\in\ZZ.
\end{equation*}
$\sigma$ 的作用可用下图表示
\begin{equation*}
\xymatrix{ x_{-3/2} \ar@/^{1pc}/[rrrr]_{\sigma}
\ar@{=} '[dr]_{C_{-1}} '[rr]_{C_{0}}  '[drrr]_{C_{1}}  [rrrr]_{C_{2}}
&  & x_{1/2}  & & x_{5/2} \\
& x_{-1/2}\ar@/_1pc/[rr]^{\sigma} &  & x_{3/2}}
\end{equation*}
图中每条双线代表一条 $\PP^{1}$。由此看出 $X_{\gamma}=\sX_{\gamma}(k)$ 只含一个点 $x_{1/2}$，即标准格 $\calO^{2}_{F}$。这意味着
\begin{equation}\label{Oga}
O_{\gamma}(\one_{\frg(\calO_{F})})=1.
\end{equation}

\subsubsection{非分歧情形：$\gamma'$} 现考虑元素 $\gamma'=\mat{0}{at^{2}}{1}{0}\in\sl_{2}(F)$。习题 \ref{exIII:stable conj} 说明 $\gamma'$ 稳定共轭于 $\gamma$，但在 $\SL_{2}(F)$ 下不与 $\gamma$ 共轭。不过 $\gamma'$ 在 $\SL_{2}(E)$ 下共轭于 $\gamma$。因此仿射 Springer 纤维 $\sX_{\gamma'}$ 在 $k_{E}$ 上仍与 $\sX_{\gamma}$ 形貌相同，只是 $\sigma\in\Gal(k_{E}/k)$ 的作用不同。

考虑 $\sX_{\gamma',k_{E}}$ 的一个分支，其 $k_{E}$-点由满足 $t\calO_{E}\oplus\calO_{E}\subset \L\subset \calO_{E}\oplus t^{-1}\calO_{E}$ 的 $\gamma'$-稳定格 $\L\subset E^{2}$ 组成。此分支由定义在 $k$ 上的条件截出，故在 $\sigma$ 下稳定，记之为 $C'_{0}$。把 $\sX_{\gamma',k_{E}}$ 的其余分支编号为 $C'_{n}$（$n\in\ZZ$），使生成元 $t\in\L_{\gamma'}$ 把 $C'_{n}$ 送到 $C'_{n+1}$。令 $x'_{n+1/2}=C'_{n}\cap C'_{n+1}\in \sX_{\gamma'}(k_{E})$。由于 $\sigma$ 在 $\sX_{\gamma,k_{E}}$ 上的作用与其在 $\L_{\gamma'}$ 上的作用（取逆）相容，唯一可能是
\begin{equation*}
\sigma(C'_{n})=C'_{-n}, \quad \sigma(x'_{n+1/2})=x'_{-n-1/2},\quad\forall n\in\ZZ.
\end{equation*}
$\sigma$ 的作用可用下图表示
\begin{equation*}
\xymatrix{ & & x'_{-1/2} \ar@/_{1pc}/@<-1ex>[r]_{\sigma}& x'_{1/2}\\
x'_{-5/2}  \ar@/_2pc/[rrrrr]^{\sigma} 
\ar@{=} '[r]^{C'_{-2}} '[urr]_{C'_{-1}}  '[urrr]_{C'_{0}}  '[rrrr]_{C'_{1}} [rrrrr]^{C'_{2}}
&  x'_{-3/2} \ar@/_{1pc}/[rrr]^{\sigma}& & & x'_{3/2}  & x'_{5/2}}
\end{equation*}
因此没有点 $x'_{n+1/2}$ 定义在 $k$ 上。分支 $C'_{0}$ 是唯一定义在 $k$ 上的分支，它在 $k$ 上必同构于 $\PP^{1}$，因为它在 $k_{E}$ 上如此（有限域上没有非平凡 Brauer--Severi 簇）。于是 $\sX_{\gamma'}(k)=C'_{0}(k)\cong\PP^{1}(k)$ 有 $q+1$ 个元素。因此
\begin{equation}\label{Ogb}
O_{\gamma'}(\one_{\frg(\calO_{F})})=q+1.
\end{equation}
把 \eqref{Oga} 与 \eqref{Ogb} 相加，得到
\begin{equation*}
SO_{\gamma}(\one_{\frg(\calO_{F})})=q+2.
\end{equation*}

\subsubsection{非分歧情形：上同调} 商 $\L_{\gamma}\bs\sX_{\gamma}$ 是结点有理曲线，由 $\PP^{1}$ 把两个 $k$-点粘成一个结点得到。

现在考虑商 $\L_{\gamma'}\bs\sX_{\gamma'}$。在 $k_{E}$ 上它同样是结点有理曲线，只有一个结点 $y$，它是所有 $x_{j}$ 的像。虽然 $y$ 是商 $\L_{\gamma'}\bs\sX_{\gamma'}$ 的 $k$-点，但它的任一原像 $x_{j}$ 都不定义在 $k$ 上。因此 $C'_{0}(k)$ 中的 $q+1$ 个点仍单射地映入商，此外还有点 $y$。于是 $(\L_{\gamma'}\bs\sX_{\gamma'})(k)$ 由 $q+2$ 个点组成。由于 $LG_{\gamma'}/\L_{\gamma'}$ 连通，$\L_{\gamma'}\bs\sX_{\gamma'}$ 的上同调的稳定部分就是整个 $\cohog{*}{\L_{\gamma'}\bs\sX_{\gamma'}}$，其上的 Frobenius 迹的交错和等于 $(\L_{\gamma'}\bs\sX_{\gamma'})(k)$ 的基数。此特例验证了公式 \eqref{d}。还要指出：$\sigma$ 在一维空间 $\cohog{1}{\L_{\gamma'}\bs\sX_{\gamma',\kbar},\Ql}$ 上的作用为 $-1$，因此 Frobenius 的 Grothendieck--Lefschetz 迹公式给出 $\#(\L_{\gamma'}\bs\sX_{\gamma'})(k)= 1+1+q$，而非 $1-1+q=q$；后者是把 $\PP^{1}$ 上两个 $k$-点等同所得的结点有理曲线上 $k$-点的个数。

\subsubsection{分歧情形：轨道积分}\label{sss:ram orb SL2} 考虑元素 $\gamma=\mat{0}{t^{2}}{t}{0}$ 与 $\gamma'=\mat{0}{at^{2}}{a^{-1}t}{0}$，其中 $a\in k^{\times}-(k^{\times})^{2}$。同样 $\gamma$ 与 $\gamma'$ 稳定共轭，但在 $\SL_{2}(F)$ 中不共轭。

设 $E$ 是 $\gamma$ 在 $\gl_{2}(F)$ 中的中心化子。则 $E$ 是 $F$ 的分歧二次扩张，$G_{\gamma}(F)=(E^{\times})^{\Nm=1}$。类似地，设 $E'$ 是 $\gamma'$ 在 $\gl_{2}(F)$ 中的中心化子。则 $E'$ 是 $F$ 的另一分歧二次扩张，$G_{\gamma'}(F)=(E'^{\times})^{\Nm=1}$。在紧群 $G_{\gamma}(F)$ 与 $G_{\gamma'}(F)$ 上选取总体积为 $1$ 的 Haar 测度。

两种情形下 $\sX_{\gamma}$ 与 $\sX_{\gamma'}$ 都同构于 $k$ 上的 $\PP^{1}$。事实上 \S\ref{SL2 ram dim 1} 已证明它们在 $\kbar$ 上同构于 $\PP^{1}$，故在 $k$ 上也必同构于 $\PP^{1}$。$\sX_{\gamma}(k)$ 与 $\sX_{\gamma'}(k)$ 都由满足 $t\calO_{F}\oplus\calO_{F}\subset \L\subset \calO_{F}\oplus t^{-1}\calO_{F}$ 的格组成，因此作为 $\Gr_{G}$ 的子簇有 $\sX_{\gamma}=\sX_{\gamma'}$。由引理 \ref{l:Ogamma}，
\begin{equation*}
O_{\gamma}(\one_{\frg(\calO_{F})})=O_{\gamma'}(\one_{\frg(\calO_{F})})=\#\sX_{\gamma}(k)=q+1.
\end{equation*}
因此
\begin{equation}\label{ram so}
SO_{\gamma}(\one_{\frg(\calO_{F})})=2(q+1).
\end{equation}

\subsubsection{分歧情形：上同调} 在 \S\ref{sss:ram orb SL2} 的设定下 $\L_{\gamma}=0$。$LG_{\gamma}$ 的分支群是 $\ZZ/2\ZZ$，但它对 $\sX_{\gamma}$ 的作用平凡。因此上同调的稳定部分就是整个 $\cohog{*}{\sX_{\gamma},\Ql}$，其上 Frobenius 的交错和给出 $\sX_{\gamma}(k)\cong\PP^{1}(k)$ 的基数。
然而 $G_{\gamma}(F)=(E^{\times})^{\Nm=1}$ 的 parahoric 子群在其中指数为 $2$（$K_{\gamma}$ 由那些在 $k$ 中约化为 $1$ 的 $e\in (E^{\times})^{\Nm=1}$ 组成）。因此公式 \eqref{d} 的右端在 $\#\sX_{\gamma}(k)$ 前多出因子 $2=\vol(K_{\gamma},d_{\gamma}g)^{-1}$。这与 \eqref{ram so} 相符，可见特例中公式 \eqref{d} 成立。

\subsection{关于基本引理的评注} 回到 \S\ref{unram sl2} 的情形。若取 $O_{\gamma'}(\one_{\frg(\calO_{F})})$ 与 $O_{\gamma}(\one_{\frg(\calO_{F})})$ 之差而非其和，会发生什么？这一差是否有类似于定理 \ref{th:orb coho} 的几何解释？

\subsubsection{$\kappa$-轨道积分} 线性组合 $O_{\gamma}(\one_{\frg(\calO_{F})})-O_{\gamma'}(\one_{\frg(\calO_{F})})$ 是 $\kappa$-轨道积分的一例。更一般地，设 $\kappa$ 是 $\cohog{1}{F,G_{\gamma}}$ 的特征标，定义 $\varphi\in \sS(\frg(F))$ 的 $\kappa$-轨道积分为
\begin{equation*}
O^{\kappa}_{\gamma}(\varphi)=\sum_{\gamma'}\kappa(\inv(\gamma,\gamma'))O_{\gamma'}(\varphi),
\end{equation*}
求和遍历 $\gamma$ 的稳定共轭类中的 $G(F)$-轨道。

\subsubsection{基本引理的陈述} Langlands--Shelstad 猜想，又称{\em 基本引理}，断言：$\gamma\in\frg(F)$ 的 $\one_{\frg(\calO_{F})}$ 的 $\kappa$-轨道积分，等于较小群 $H$ 的某个元素 $\gamma_{H}\in\frh(F)$ 的稳定轨道积分，相差一个简单因子。用公式写出，基本引理即恒等式
\begin{equation*}
O^{\kappa}_{\gamma}(\one_{\frg(\calO_{F})})=\Delta(\gamma,\gamma_{H})\cdot SO_{\gamma_{H}}(\one_{\frh(\calO_{F})}).
\end{equation*}
较小群 $H$ 同时依赖于 $G$ 与 $\kappa$，称为 $(G,\kappa)$ 的{\em 内窥群}。数 $\Delta(\gamma,\gamma_{H})$ 称为{\em 转移因子}；当 $\gamma$ 从它的稳定共轭类中适当选取、且 $G_{\gamma}$ 与 $H_{\gamma_{H}}$ 上的测度选取适当时，它是 $q$ 的整数次幂（依赖于 $\gamma$ 与 $\gamma_{H}$）。

\subsubsection{一个简单情形} 在 \S\ref{unram sl2} 的情形中，取 $\cohog{1}{F,G_{\gamma}}\cong\ZZ/2\ZZ$ 上的非平凡特征标 $\kappa$，相应的内窥群 $H$ 同构于环面 $G_{\gamma}$；但一般情形下它未必同构于 $G$ 的子群。此时基本引理是恒等式
\begin{equation*}
O_{\gamma'}(\one_{\frg(\calO_{F})})-O_{\gamma}(\one_{\frg(\calO_{F})})=q=q\cdot SO_{\gamma_{H}}(\one_{\frh(\calO_{F})}),
\end{equation*}
其中把 $H$ 与 $G_{\gamma}$ 等同后取 $\gamma_{H}=\gamma$。

另一方面，在 \S\ref{sss:ram orb SL2} 的分歧情形中，$\gamma$ 关于非平凡 $\kappa$ 的 $\kappa$-轨道积分消没。这也许无需显式计算轨道积分即可解释：一般说来，为使 $O^{\kappa}_{\gamma}(\one_{\frg(\calO_{F})})$ 可能非零，$\kappa$ 必须经过 $\cohog{1}{F,G_{\gamma}}$ 的进一步商群分解。确切陈述见 \cite[命题 8.2.7]{NgoFL}。

\subsubsection{关于证明的评注} 一般 $G$ 与函数域 $F$ 的基本引理由 B-C. Ng\^o 证明 \cite{NgoFL}。定理 \ref{th:orb coho} 可以推广到 $\kappa$-轨道积分，只需把 $\L_{\gamma}\bs\sX_{\gamma}$ 的上同调的稳定部分换成 $\kappa$-同型部分。Goresky、Kottwitz 与 MacPherson \cite{GKM} 用这一推广把基本引理重述为仿射 Springer 纤维 $\sX_{\gamma}$ 与其内窥表亲 $\sX^{H}_{\gamma_{H}}$ 的上同调群之间的关系。他们还证明了一些特殊但极不平凡情形下的基本引理。Ng\^o 的证明建立在这一上同调重述之上，但另用了一个新要素，即 Hitchin 纤维——它可以看作仿射 Springer 纤维的「整体」类比。这将是下一讲的主题。

\subsection{习题}

本组习题中，$k=\FF_{q}$ 是 $\chark\neq2$ 的有限域，$F=k\lr{t}$。

\subsubsection{}\label{exIII:K inv} 设 $H$ 是 $F$ 上代数群，$f\in\sS(H(F))$。证明存在紧开子群 $K\subset H(F)$，使 $f$ 在 $K$ 下左右不变。

\subsubsection{}\label{exIII:Ogamma} 证明引理 \ref{l:Ogamma}。

\subsubsection{}\label{exIII:stable conj} 设 $a\in k^{\times}-(k^{\times})^{2}$，$G=\SL_{2}$。考虑 \S\ref{sss:st conj SL2} 中的矩阵 $\gamma,\gamma'$。证明 $\gamma$ 与 $\gamma'$ 稳定共轭，但在 $\SL_{2}(F)$ 下不共轭。

\subsubsection{} 设 $F$ 上 $G=\GL_{n}$，$\gamma=\diag(\gamma_{1},\cdots,\gamma_{m})$ 是 $\frg(F)^{\rs}$ 中的块对角矩阵，其中 $\gamma_{i}\in\mathfrak{gl}_{n_{i}}(F)^{\rs}$。考虑轨道积分 $O_{\gamma}(\one_{\frg(\calO_{F})})$ 当 $q=\#k$ 趋于无穷时的渐近行为。求最小的整数 $d$，使
\begin{equation*}
O_{\gamma}(\one_{\frg(\calO_{F})})=O(q^{d})
\end{equation*}
（$q\to\infty$）。注意右端的 $O(\cdot)$ 是分析学家的 $O$，而左端的 $O$ 指轨道积分\zy{原文此处玩双关：$O_{\gamma}$ 的 $O$ 是 orbital integral 的首字母，$O(q^{d})$ 的 $O$ 是大 $O$ 记号。译文保留两者同形。来源：原文 \texttt{ZhiweiYunPCMI.tex} 第 2141 行原文的自注。}。能否用 $\gamma_{i}$ 的特征多项式解释 $d$？

关于 $O_{\gamma}(\one_{\frg(\calO_{F})})$ 的显式估计，见 \cite{Y-Orb}。

\subsubsection{} 设 $G=\SL_{3}$，$\gamma=\left(\begin{array}{ccc} 0 & 0 & t^{4}\\ 1 & 0 & 0 \\ 0 & 1& 0 \end{array}\right)\in\frg(F)$。计算 $O_{\gamma}(\one_{\frg(\calO_{F})})$。

{\em 提示}：使用 \S\ref{sss:LS} 中引入的胞腔分解。

\subsubsection{} 设 $G=\SL_{2}$。设 $f$ 是满足以下条件的元素 $X\in\frg(\calO_{F})$ 的特征函数：其约化 $\ov{X}\in\frg(k)$ 是正则幂零元。设 $\gamma\in\frg(F)$ 是正则半单元。
\begin{enumerate}
\item 证明除非 $\det(\gamma)\in t\calO_{F}$，有 $O_{\gamma}(f)=0$。
\item 当 $\det(\gamma)\in t\calO_{F}$ 时，证明
\begin{equation*}
O_{\gamma}(f)=O_{\gamma}(\one_{\frg(\calO_{F})})-O_{t^{-1}\gamma}(\one_{\frg(\calO_{F})}).
\end{equation*}
\end{enumerate}

\subsubsection{} 设 $G=\GL_{2}$，$\gamma=\mat{0}{t^{n}}{1}{0}$，其中 $n\geq 1$ 为奇数。设 $G(F)_{d}$ 是满足 $\val_{F}(\det g)=d$ 的 $g\in G(F)$ 之集。固定 $G(F)$ 上的 Haar 测度，使 $G(\calO_{F})$ 的体积为 $1$。证明对任意整数 $d\geq0$，
\begin{equation*}
\int_{G(F)_{d}}\one_{\frg(\calO_{F})}(g^{-1}\gamma g)\one_{\calO^{2}_{F}}((0,1)g)dg
\end{equation*}
等于平面曲线 $y^{2}-t^{n}=0$ 的闭子概形 $Z$ 的个数，其中 $Z$ 满足：(1) $Z$ 的底拓扑空间是点 $(y,t)=(0,0)$；(2) $\dim_{k}\calO_{Z}=d$。

注：此习题把轨道积分与曲线上的点 Hilbert 概形联系起来。这一联系已被 \cite{Y-Orb} 用来给出 $\GL_{n}$ 的轨道积分的估计。整体类比见 \S\ref{sss:MHilb} 与 \S\ref{sss:Hilb}。
